% !PRINTED 433
\weberpage{433}
%<<CONT
oder
\weq{14}{\beta\,\omega^2 + (\beta\,c'+\beta\,c+\alpha\,p-\delta\,p)\,\omega + \beta\,c\,c'+\alpha\,c\,p-\delta\,c'\,p-\gamma\,p^2 = 0,}
woraus durch Vergleichung mit (4):
\weq{15}{\begin{aligned}
\beta &= A\,x,\\
\beta\,c'+\beta\,c+\alpha\,p-\delta\,p &= B\,x,\\
\beta\,c\,c'+\alpha\,c\,p-\delta\,c'\,p-\gamma\,p^2 &= C\,x.
\end{aligned}}
und wenn man wieder
\weq{16}{\beta\,c'-\beta\,c+\alpha\,p+\delta\,p = y}
setzt:
\weq{17}{\begin{aligned}
2\,(\beta\,c'+\alpha\,p) &= B\,x+y,\\
2\,(\beta\,c-\delta\,p) &= B\,x-y.
\end{aligned}}
Daraus
\weq{18}{y^2+m\,x^2 = 4\,p^2.}
Ist $x$ durch $p$ teilbar, so muß auch $y$ durch $p$ teilbar sein, und wir kommen durch Wegheben des Faktors $p^2$ auf die Gleichung (11) zurück, die, wie wir gesehen haben, nur für die Fälle $m=3$ und $m=4$ lösbar ist. Wir nehmen also weiter an, $x$ sei durch $p$ unteilbar.

Ist $m$ durch $p$ teilbar, so ist auch $y$ durch $p$ teilbar, $\beta$ ist nach (15) nicht durch $p$ teilbar. Dann aber folgt aus (16): $c\equiv c' \pmod{p}$, und beide entsprechen also der nämlichen Wurzel von (1).

Es gehe jetzt also $p$ weder in $x$ noch in $m$ auf. Dann geht es auch nicht in $y$ auf, und nach (17) werden $c$ und $c'$ aus den Kongruenzen
\weq{19}{\beta\,c' \equiv \frac{B\,x+y}{2}, \quad \beta\,c \equiv \frac{B\,x-y}{2} \pmod{p}}
bestimmt, und wenn $c,\,c'$ bestimmt sind, ergeben sich aus (17) $\alpha$ und $\delta$ als \gesperrt{ganze Zahlen}. Damit sind die beiden ersten Gleichungen (15) befriedigt, und $\gamma$ erhält man aus der letzten Gleichung (15) gleichfalls als ganze Zahl.

Es ist nämlich
\weq{20}{\gamma\,p^2 = \beta\,c\,c'+(\alpha\,c-\delta\,c')\,p-C\,x \equiv 0 \pmod{p^2}.}
Denn wir haben aus (17) und (15)
\[
\begin{aligned}
\beta^2\,c\,c'+\beta\,(\alpha\,c-\delta\,c')\,p-\alpha\,\delta\,p^2-A\,C\,x^2
&= \beta\,[\beta\,c\,c'+(\alpha\,c-\delta\,c')\,p-C\,x]-\alpha\,\delta\,p^2\\
&= -\frac{y^2+m\,x^2}{4} = -p^2;
\end{aligned}
\]
%CONT>>
